\section{Experiments}
\label{sec:experiments}

\subsection{Datasets, metrics, and splits}

We evaluate on six standard change-detection benchmarks:

\begin{itemize}
  \item \textbf{LEVIR-CD}~\citep{chen2020spatial} — building change, 7120/1024/2048
    train/val/test 256$^2$ crops of the standard split.
  \item \textbf{LEVIR-CD+}~\citep{zheng2022changestar} — extended LEVIR, 637/348
    train/test 1024$^2$ tiles; we hold out 10\% of train as val.
  \item \textbf{S2Looking}~\citep{shen2021s2looking} — 3500/500/1000 rural,
    off-nadir 1024$^2$ tiles.
  \item \textbf{CDD}~\citep{lebedev2018change} — 10000/2998/3000 256$^2$ crops
    with seasonal variation.
  \item \textbf{DSIFN-CD}~\citep{zhang2020deeply} — we use the 10000/2998/3000
    split from the official \texttt{Real/subset}; we name the split explicitly
    because cross-paper comparisons can differ by $\sim$5 F1.
  \item \textbf{SECOND}~\citep{yang2021second} (binary collapse) —
    2968/1694 train/test; val is a 10\% auto-split. Binary change masks are
    generated from the semantic labels as $M = \mathbf{1}[\ell^{(r)} \ne \ell^{(t)}]$.
\end{itemize}

We report F1, IoU, Precision, Recall, Overall Accuracy (OA), and Cohen's
$\kappa$ using pixel-level TP/FP/FN/TN accumulated at threshold 0.5.
On CDD, F1 saturates above 95 in the recent literature; we follow the
community practice~\citep{corley2024reality} of reporting IoU as the primary
discriminator.

\subsection{Main results}

Table~\ref{tab:sota} compares GALASH to published SOTA on each benchmark.
Our numbers are the TODO.

\begin{table}[t]
  \centering
  \small
  \caption{Per-benchmark comparison to published SOTA (F1 / IoU).
    Bold: best. Underline: second-best. All GALASH numbers use
    \texttt{dinov2\_rs\_base + sam2\_base\_plus} with the frozen decoder unless
    noted. TTA denotes 8-way test-time augmentation.}
  \label{tab:sota}
  \begin{tabular}{lcccccc}
    \toprule
    Method & LEVIR-CD & LEVIR-CD+ & S2Looking & CDD & DSIFN-CD & SECOND \\
    \midrule
    FC-Siam-diff~\citep{daudt2018fully}    & 86.31 & ---    & 13.1  & 70.61 & 70.55 & ---   \\
    STANet~\citep{chen2020spatial}         & 87.26 & ---    & ---   & 84.12 & 64.56 & ---   \\
    BIT~\citep{chen2021remote}             & 89.31 & ---    & 61.85 & 88.90 & 87.61 & ---   \\
    SNUNet~\citep{fang2021snunet}          & 88.16 & ---    & ---   & 83.89 & 66.18 & ---   \\
    ChangeFormer~\citep{bandara2022changeformer} & 91.11 & ---    & 63.39 & 94.63 & 94.67 & ---   \\
    TinyCD~\citep{codegoni2023tinycd}      & 91.05 & ---    & ---   & ---   & ---   & ---   \\
    ChangeStar~\citep{zheng2022changestar} & ---   & 90.5   & ---   & ---   & ---   & ---   \\
    ChangerEx~\citep{fang2023changer}      & 91.77 & ---    & 66.20 & ---   & ---   & ---   \\
    DDPM-CD~\citep{bandara2024ddpmcd}      & 90.91 & ---    & ---   & 95.62 & \textbf{96.65} & ---   \\
    BAN~\citep{li2024ban}                  & 91.86 & ---    & 66.70 & ---   & ---   & ---   \\
    TTP~\citep{chen2024ttp}                & 92.10 & ---    & ---   & ---   & ---   & ---   \\
    SChanger~\citep{sun2025schanger}       & \textbf{92.87} & ---    & \underline{68.95} & \textbf{97.62} & ---   & ---   \\
    UniChange~\citep{zeng2025unichange}    & ---   & ---    & \textbf{69.32} & ---   & ---   & \textbf{73.12} \\
    \midrule
    GALASH (ours, frozen)                  & TODO  & TODO   & TODO  & TODO  & TODO  & TODO  \\
    GALASH (ours, fine-tuned)              & TODO  & TODO   & TODO  & TODO  & TODO  & TODO  \\
    \bottomrule
  \end{tabular}
\end{table}

\subsection{Foundation-model sweep}

A central question of this paper is how the choice of encoder and decoder
families interacts. Table~\ref{tab:enc-dec} reports F1 on LEVIR-CD across
a representative subset of encoder $\times$ decoder pairings.

\begin{table}[t]
  \centering
  \small
  \caption{Encoder $\times$ decoder sweep on LEVIR-CD (F1, frozen decoder).}
  \label{tab:enc-dec}
  \begin{tabular}{lccccc}
    \toprule
    \multirow{2}{*}{Encoder} & \multicolumn{4}{c}{Decoder} & trainable \\
    \cmidrule(lr){2-5}
    & sam2-tiny & sam2-small & sam2-base+ & sam2-large & params (M) \\
    \midrule
    DINOv2-base               & TODO & TODO & TODO & TODO & 10.5 \\
    DINOv2-RS-base            & TODO & TODO & TODO & TODO & 10.5 \\
    DINOv2-RS-large           & TODO & TODO & TODO & TODO & 18.2 \\
    DINOv3-base               & TODO & TODO & TODO & TODO & 10.5 \\
    DINOv3-large              & TODO & TODO & TODO & TODO & 18.2 \\
    DINOv3-sat-large          & TODO & TODO & TODO & TODO & 18.2 \\
    \bottomrule
  \end{tabular}
\end{table}

\subsection{Cross-dataset generalisation}

To test whether GALASH's frozen-foundation design generalises beyond the
training distribution, Table~\ref{tab:xdataset} reports F1 obtained when
training on one benchmark and evaluating on another. Values along the diagonal
are the in-distribution results from Table~\ref{tab:sota}.

\begin{table}[t]
  \centering
  \small
  \caption{Cross-dataset generalisation F1.
    Rows: training source. Columns: evaluation target.}
  \label{tab:xdataset}
  \begin{tabular}{lcccccc}
    \toprule
    Train $\downarrow$ / Test $\rightarrow$ &
    LEVIR-CD & LEVIR-CD+ & S2Looking & CDD & DSIFN-CD & SECOND \\
    \midrule
    LEVIR-CD          & TODO & TODO & TODO & TODO & TODO & TODO \\
    LEVIR-CD+         & TODO & TODO & TODO & TODO & TODO & TODO \\
    S2Looking         & TODO & TODO & TODO & TODO & TODO & TODO \\
    CDD               & TODO & TODO & TODO & TODO & TODO & TODO \\
    DSIFN-CD          & TODO & TODO & TODO & TODO & TODO & TODO \\
    SECOND            & TODO & TODO & TODO & TODO & TODO & TODO \\
    Pooled            & TODO & TODO & TODO & TODO & TODO & TODO \\
    \bottomrule
  \end{tabular}
\end{table}
